\section*{Bab \ref{ch:g10:the-mole} --- Mol dan Massa Molar}

\begin{solution}{exo:g10:the-mole:1}
$M(\ce{O2}) = 2 \times 16.0 = \qty{32.0}{g/mol}$;
$M(\ce{N2}) = 2 \times 14.0 = \qty{28.0}{g/mol}$;
$M(\ce{CH4}) = 12.0 + 4 \times 1.0 = \qty{16.0}{g/mol}$;
$M(\ce{NH3}) = 14.0 + 3 \times 1.0 = \qty{17.0}{g/mol}$;
$M(\ce{CO2}) = 12.0 + 2 \times 16.0 = \qty{44.0}{g/mol}$.
\end{solution}

\begin{solution}{exo:g10:the-mole:2}
Air: $n = 36.0 / 18.0 = \qty{2.00}{mol}$. Metana:
$n = 4.0 / 16.0 = \qty{0.25}{mol}$.
\end{solution}

\begin{solution}{exo:g10:the-mole:3}
$m(\ce{NaCl}) = 0.25 \times 58.5 = \qty{14.6}{g}$, sekitar \qty{15}{g};
$m(\ce{CO2}) = 2.0 \times 44.0 = \qty{88}{g}$.
\end{solution}

\begin{solution}{exo:g10:the-mole:4}
$N = 0.50 \times \num{6.02e23} = \num{3.0e23}$ \omterm{def:g7:atoms-and-molecules:molecule}{molekul} gas oksigen.
$n = \num{3.0e24} / \qty{6.02e23}{mol^{-1}} = \qty{5.0}{mol}$ besi.
\end{solution}

\begin{solution}{exo:g10:the-mole:5}
$V = n \times V_m = 0.50 \times 24.1 = \qty{12}{L}$.
$n = V / V_m = 12 / 24.1 = \qty{0.50}{mol}$.
\end{solution}

\begin{solution}{exo:g10:the-mole:6}
Tetes itu bermassa \qty{0.050}{g}, jadi
$n = 0.050 / 18.0 = \qty{2.8e-3}{mol}$ dan
$N = \num{2.8e-3} \times \num{6.02e23} = \num{1.7e21}$ \omterm{def:g7:atoms-and-molecules:molecule}{molekul}.
\end{solution}

\begin{solution}{exo:g10:the-mole:7}
Sebuah \omterm{def:g7:atoms-and-molecules:atom}{atom} aluminium (\qty{27.0}{g/mol}) kira-kira setengah berat \omterm{def:g7:atoms-and-molecules:atom}{atom}
besi (\qty{55.8}{g/mol}), sehingga aluminium bermassa sama memuat
kira-kira dua kali lebih banyak \omterm{def:g7:atoms-and-molecules:atom}{atom}. Periksa: aluminium
$10/27.0 = \qty{0.370}{mol}$, yaitu \num{2.23e23} \omterm{def:g7:atoms-and-molecules:atom}{atom}; besi
$10/55.8 = \qty{0.179}{mol}$, yaitu \num{1.08e23} \omterm{def:g7:atoms-and-molecules:atom}{atom}. Aluminium
menang, dengan faktor sekitar 2.1.
\end{solution}

\begin{solution}{exo:g10:the-mole:8}
Satu liter adalah $1/24.1 = \qty{0.0415}{mol}$ gas. Karbon dioksida:
$0.0415 \times 44.0 = \qty{1.83}{g}$; gas nitrogen:
$0.0415 \times 28.0 = \qty{1.16}{g}$. Karbon dioksida kira-kira satu
setengah kali lebih rapat daripada \omterm{def:g4:air-a-mixture-of-gases:air}{udara}, yang \omterm{def:g10:the-mole:molar-mass}{massa molarnya} dekat
dengan \omterm{def:g10:the-mole:molar-mass}{massa molar} gas nitrogen. Sari anggur yang berfermentasi
melepaskan karbon dioksida, yang turun dan menumpuk di dekat lantai
gudang bawah tanah tertutup, tempat gas itu dapat membuat sesak napas
siapa saja yang masuk.
\end{solution}

\begin{solution}{exo:g10:the-mole:9}
$M(\ce{C6H12O6}) = 6 \times 12.0 + 12 \times 1.0 + 6 \times 16.0 =
\qty{180.0}{g/mol}$. Anak panah $\div N_A$: $n = \num{1.0e22} / \num{6.02e23} =
\qty{1.66e-2}{mol}$; anak panah $\times M$: $m = \num{1.66e-2} \times 180.0 =
\qty{3.0}{g}$.
\end{solution}

\begin{solution}{exo:g10:the-mole:10}
$M = \qty{342.0}{g/mol}$, jadi $n = 5.0 / 342.0 = \qty{1.46e-2}{mol}$;
$N = \num{1.46e-2} \times \num{6.02e23} = \num{8.8e21}$ \omterm{def:g7:atoms-and-molecules:molecule}{molekul}. Setiap
\omterm{def:g7:atoms-and-molecules:molecule}{molekul} memiliki 12 \omterm{def:g7:atoms-and-molecules:atom}{atom} karbon: $12 \times \num{8.8e21} = \num{1.1e23}$
\omterm{def:g7:atoms-and-molecules:atom}{atom} karbon.
\end{solution}

\begin{solution}{exo:g10:the-mole:11}
Air: $1000 / 18.0 = \qty{55.6}{mol}$. Etanol:
$M = 2 \times 12.0 + 6 \times 1.0 + 16.0 = \qty{46.0}{g/mol}$, jadi
$1000 / 46.0 = \qty{21.7}{mol}$. Air memiliki \omterm{def:g10:the-mole:amount}{jumlah zat} yang lebih
besar, dengan faktor $55.6 / 21.7 \approx 2.6$, yang tidak lain adalah
$46.0 / 18.0$.
\end{solution}

\begin{solution}{exo:g10:the-mole:12}
Emas: $\qty{3.75}{g}$, $n = 3.75 / 197.0 = \qty{1.90e-2}{mol}$, yaitu
\num{1.15e22} \omterm{def:g7:atoms-and-molecules:atom}{atom}. Tembaga: \qty{1.25}{g},
$n = 1.25 / 63.5 = \qty{1.97e-2}{mol}$, yaitu \num{1.19e22} \omterm{def:g7:atoms-and-molecules:atom}{atom}.
Ternyata cincin itu memuat \omterm{def:g7:atoms-and-molecules:atom}{atom} tembaga sedikit lebih banyak daripada
\omterm{def:g7:atoms-and-molecules:atom}{atom} emas: sebuah \omterm{def:g7:atoms-and-molecules:atom}{atom} emas kira-kira tiga kali lebih berat daripada
sebuah \omterm{def:g7:atoms-and-molecules:atom}{atom} tembaga.
\end{solution}

\begin{solution}{exo:g10:the-mole:13}
\begin{enumerate}
  \item $V = 8.0 \times 6.0 \times 3.0 = \qty{144}{m^3} =
        \qty{1.44e5}{L}$; $n = \num{1.44e5} / 24.1 = \qty{5.98e3}{mol}$;
        $N = \num{5.98e3} \times \num{6.02e23} = \num{3.6e27}$ \omterm{def:g7:atoms-and-molecules:molecule}{molekul}.
  \item $M = (78 \times 28.0 + 21 \times 32.0 + 1 \times 40.0) / 100 =
        2896 / 100 = \qty{29.0}{g/mol}$.
  \item $m = \num{5.98e3} \times 29.0 = \qty{1.73e5}{g}$, sekitar
        \qty{173}{kg}: \omterm{def:g4:air-a-mixture-of-gases:air}{udara} di sebuah ruang kelas sama beratnya dengan
        dua atau tiga orang dewasa.
\end{enumerate}
\end{solution}

\begin{solution}{exo:g10:the-mole:14}
\begin{enumerate}
  \item $n = 1000 / 28.1 = \qty{35.6}{mol}$, jadi
        $N = 35.6 \times \num{6.02e23} = \num{2.14e25}$ \omterm{def:g7:atoms-and-molecules:atom}{atom}.
  \item $m = \qty{1.000}{kg} / \num{2.14e25} = \qty{4.67e-26}{kg}$.
  \item Massa bola dan \omterm{def:g10:the-mole:molar-mass}{massa molar} memberikan \omterm{def:g10:the-mole:amount}{jumlah zatnya},
        $n = m / M$. Jika banyaknya \omterm{def:g7:atoms-and-molecules:atom}{atom} $N$ dihitung secara terpisah
        (dari volume bola dan jarak antaratom dalam kristal), rasio
        $N / n$ adalah \omterm{def:g10:the-mole:avogadro}{tetapan Avogadro}.
\end{enumerate}
\end{solution}

\begin{solution}{exo:g10:the-mole:15}
$M = 0.758 \times 35.0 + 0.242 \times 37.0 = 26.53 + 8.95 =
\qty{35.5}{g/mol}$: nilai pada tabel.
\end{solution}

\begin{solution}{pb:g10:the-mole:1}
\textbf{1.} $M(\ce{H2O}) = 2 \times 1.0 + 16.0 = \qty{18.0}{g/mol}$.

\textbf{2.} $m = 250 \times \qty{1.0}{g} = \qty{250}{g}$.

\textbf{3.} $n = 250 / 18.0 = \qty{13.9}{mol}$.

\textbf{4.} Setiap \omterm{def:g7:atoms-and-molecules:molecule}{molekul} memiliki dua \omterm{def:g7:atoms-and-molecules:atom}{atom} hidrogen dan satu \omterm{def:g7:atoms-and-molecules:atom}{atom}
oksigen: \qty{27.8}{mol} \omterm{def:g7:atoms-and-molecules:atom}{atom} hidrogen, \qty{13.9}{mol} \omterm{def:g7:atoms-and-molecules:atom}{atom} oksigen.

\textbf{5.} $N = 13.9 \times \num{6.02e23} = \num{8.36e24}$ \omterm{def:g7:atoms-and-molecules:molecule}{molekul}.

\textbf{6.} $m = \qty{18.0}{g/mol} / \qty{6.02e23}{mol^{-1}} =
\qty{2.99e-23}{g}$.

\textbf{7.} $\num{8.36e24} \times \qty{2.99e-23}{g} = \qty{250}{g}$,
yaitu massa air di dalam gelas, sebagaimana mestinya.

\textbf{8.} $\num{8.36e24} / \num{3.16e7} = \num{2.6e17}$ tahun: bilangan
tahun dengan tujuh belas nol, jauh melampaui hitungan yang mungkin.

\textbf{9.} $\qty{1.335e9}{km^3} \times \num{e9} = \qty{1.335e18}{m^3}$,
dan dengan $\qty{1}{m^3} = \qty{1000}{L}$, \qty{1.335e21}{L}.

\textbf{10.} $\num{8.36e24} / \num{1.335e21} = \num{6.26e3}$ \omterm{def:g7:atoms-and-molecules:molecule}{molekul}
yang ditandai per liter.

\textbf{11.} $\num{6.26e3} \times 0.250 = \num{1.57e3}$ \omterm{def:g7:atoms-and-molecules:molecule}{molekul} yang
ditandai dalam gelas kedua.

\textbf{12.} Massa $\num{1.335e21} \times \qty{1.0}{kg} =
\qty{1.335e21}{kg} = \qty{1.335e24}{g}$; $n = \num{1.335e24} / 18.0 =
\qty{7.42e22}{mol}$.

\textbf{13.} $13.9 / \num{7.42e22} = \num{1.87e-22}$.

\textbf{14.} $\num{1.87e-22} \times \num{8.36e24} = \num{1.57e3}$:
jawaban yang sama dengan soal 11, sebagaimana mestinya.

\textbf{15.} $\num{1.335e21} / 0.250 = \num{5.34e21}$ gelas.

\textbf{16.} $\num{8.36e24} / \num{5.34e21} \approx 1570$. Satu gelas
memuat kira-kira 1600 kali lebih banyak \omterm{def:g7:atoms-and-molecules:molecule}{molekul} daripada banyaknya gelas
air di samudra; jika tersebar merata, \omterm{def:g7:atoms-and-molecules:molecule}{molekul-molekul} dari satu gelas
menyisakan sekitar 1600 di setiap gelas air samudra.

\textbf{17.} $1 / \num{6.26e3} = \qty{1.6e-4}{L} = \qty{0.16}{mL}$,
kira-kira tiga tetes.

\textbf{18.} Volume samudra adalah perkiraan yang paling-paling
diketahui sampai tiga atau empat angka, dan telah dibulatkan; air laut
bukan air murni (air laut mengandung garam); gelasnya tidak tepat
\qty{250}{mL}; dan \omterm{def:g2:mixing-and-dissolving:mix}{pencampuran} ke seluruh samudra tidak akan pernah
benar-benar merata. Hanya orde besarnya yang dapat dipercaya.

\textbf{19.} Sekitar 1600 \omterm{def:g7:atoms-and-molecules:molecule}{molekul} dari gelas pertama kembali dalam gelas
kedua.
\end{solution}
